Every result in this subsection runs the identical, unmodified pipeline ---
the same conformal calibration, the same episodic e-process, the same
classical baselines --- on real per-frame scores rather than the simulator.
Two classical video-anomaly-detection corpora supply continuous, real,
frame-labelled surveillance footage: CUHK Avenue \citep{lu2013avenue}
(\RealAvenueNTrainVideos{} training / \RealAvenueNTestVideos{} testing
videos, a campus walkway, \RealAvenueNEpisodes{} distinct labelled anomalous
runs across the test split) and UCSD Ped2 \citep{mahadevan2010ucsd}
(\RealPedTwoNTrainVideos{} training / \RealPedTwoNTestVideos{} testing videos, a
pedestrian walkway, \RealPedTwoNEpisodes{} runs, one per test video). Neither
ships per-frame detector scores, so we compute our own: Farneback dense
optical-flow magnitude \citep{farneback2003}, averaged per frame and denoised
by a 3-frame per-video median filter (disclosed below), reaching frame-level
ROC-AUC \RealPedTwoAuc{} on Ped2 and \RealAvenueAuc{} on Avenue --- a real,
classical, unsupervised backbone chosen for transparency over benchmark
leaderboard position; \cref{sec:related} situates it against 2025--2026
methods that reach substantially higher AUC on the same pair. Streams are
built by \texttt{build\_stream\_from\_clips} exactly as \cref{sec:problem}
describes: calibration is the official training split, and the deployment
stream is spliced from real labelled runs inside the test split, with no
frame ever eligible for both roles.

\paragraph{A real train/test shift, disclosed and quantified.} Before any
sequential machinery runs, a direct quantile check on the raw scores shows
that \RealPedTwoShiftFrac{}\% of Ped2's confirmed-normal test frames, and
\RealAvenueShiftFrac{}\% of Avenue's, exceed the 95th percentile of the
\emph{official training} split's score distribution --- both far above the
5\% a well-calibrated split would show. Real footage across a train/test
boundary is not, in this backbone's scores, the same distribution twice; this
is a fact about the datasets and the backbone, not an artefact of splicing.

\paragraph{Real autocorrelation, and why it is the headline finding here.}
\Cref{thm:beta} needs the calibration draws to be close to i.i.d., which is
why \cref{sec:layer1} thins the betting grid to the decorrelation lag. On the
AR(1) simulator that lag is small and an anomalous episode lasts thousands of
frames, so thinning costs almost nothing. Real Farneback-flow scores behave
nothing like that: the sample autocorrelation of the official training split
does not fall below 0.05 until lag \RealPedTwoFullLag{} on Ped2 and lag
\RealAvenueFullLag{} on Avenue --- two orders of magnitude past anything the
simulator produced --- while a real labelled episode lasts a median of
\RealPedTwoEpiMedian{} frames on Ped2 (\RealPedTwoEpiMin{}--\RealPedTwoEpiMax{}) and
just \RealAvenueEpiMedian{} on Avenue (\RealAvenueEpiMin{}--\RealAvenueEpiMax{}).
Thinning to the fully rigorous lag leaves at most one or two betting
opportunities inside a typical real event --- too few for any sequential test,
at any $\alpha$, to detect reliably.

\begin{figure}[t]
  \centering
  \resultfig{fig_r1_validity}{}
  \caption{Real video: validity and power against the betting-grid lag
  (\RealLagReps{} streams per point, $\alpha=0.1$), and time-uniform
  false-alarm control at the fixed operating lag each dataset's sweep
  selects. \textbf{(a)} Below a sharp threshold the realised false-alarm rate
  (solid) is uncontrolled; above it the miss rate (dashed) is near one. The
  two curves cross in a window of a handful of frames --- not the wide margin
  the simulated streams left --- and the vertical line marks the lag used
  everywhere else in this subsection. \textbf{(b)} PFA against the nominal
  level at that fixed lag; open markers are zero-count outcomes plotted at
  their 95\% Wilson upper limit.}
  \label{fig:real-validity}
\end{figure}

\Cref{tab:real_lag_sweep} reports that crossing directly. Rather than pick
the one lag that works best and report only it, every lag swept is shown,
because the crossing itself --- not the single surviving operating point ---
is the finding: real per-frame anomaly scores do not leave a comfortable
region in which both validity and power hold simultaneously, the way the
simulated AR(1) streams did throughout \cref{sec:validity,sec:episodes}. The
lag this sweep selects (\RealPedTwoLag{} on Ped2, \RealAvenueLag{} on Avenue)
is fixed once, from this diagnostic alone, and reused unchanged for every
other real-data result in this subsection --- including the classical
baselines in \cref{tab:real_delay_far}, which are thinned to the identical
grid so that no method is favoured by a different frame rate.

\inputtable{tabr1_lag_sweep}

\Cref{tab:real_validity} sweeps $\alpha$ at that fixed lag.
Ped2 controls its false-alarm rate at \RealPedTwoPfaHeadline{}
(95\% CI upper end \RealPedTwoCiHiHeadline{}) at $\alpha=0.1$, with roughly two
thirds of real spliced events detected (miss rate \RealPedTwoMissHeadline{}) --- a
real, modest, but genuine detection window. Avenue is harder: its shorter,
sparser real episodes leave a realised PFA of \RealAvenuePfaHeadline{} at the same
$\alpha$, but a miss rate of \RealAvenueMissHeadline{}. At the loosest level swept,
$\alpha=\RealAvenueWorstAlpha$, Avenue's realised rate (\RealAvenueWorstPfa{})
still does not exceed the nominal level, so the guarantee itself is not
violated anywhere on this grid; what real footage costs Avenue is power, not
validity. This is not the failure the simulated experiments were built to
surface, and it could not have been --- it is specific to what real
optical-flow scores from real short-duration events actually look like, and
\cref{sec:limitations} returns to it as the paper's central open problem
rather than a number to look past.

\inputtable{tabr1b_validity}

\paragraph{Classical baselines on the same real streams.}
\Cref{tab:real_delay_far} repeats the delay-vs-false-alarm-rate comparison of
\cref{sec:delay} on real video, with every method --- SENTINEL-E, a fixed
threshold, the per-frame p-value threshold, CUSUM, Shiryaev--Roberts, a
parametric e-detector, and E-SHIFT --- thinned to the identical fixed lag
\cref{tab:real_lag_sweep} selects for that dataset. SENTINEL-E is the only
method that reaches the target false-alarm probability while still detecting
most events on either dataset (\RealDelayPedTwoNBaselinesPass{} of
\RealDelayPedTwoNBaselines{} classical baselines pass on Ped2,
\RealDelayAvenueNBaselinesPass{} of \RealDelayAvenueNBaselines{} on Avenue).
The per-frame p-value threshold --- consuming the identical conformal
p-values SENTINEL-E does, differing only in the stopping rule --- never fires
at any swept threshold on either dataset: no single real frame's score is
extreme enough on its own to cross even a lenient per-frame significance
level, which is exactly the argument \cref{sec:episodic} makes for accumulating
evidence sequentially rather than thresholding one frame at a time, now
confirmed on real scores rather than simulated ones. The classical
parametric procedures (CUSUM, Shiryaev--Roberts, the parametric e-detector)
reach a false-alarm probability near one at essentially every setting of
their own tuning knob on real footage, consistent with \cref{sec:delay}'s
simulated finding and for the same underlying reason: their pre-change model
is a poor fit to real, heavy-tailed, strongly autocorrelated detector output.

\begin{figure}[t]
  \centering
  \resultfig{fig_r2_delay_far}{}
  \caption{Detection delay against realised false-alarm probability on real
  video (\RealDelayReps{} null and signal streams per operating point), at
  each dataset's fixed operating lag from \cref{fig:real-validity}. Curves
  that never reach the plotted range detect too rarely, on real spliced
  events, at any realised false-alarm rate their swept knob achieves.}
  \label{fig:real-delayfar}
\end{figure}

\inputtable{tabr2_delay_far}

\paragraph{What this subsection does and does not establish.} It establishes
that the pipeline runs end to end on real, unmodified video --- real
backbone, real train/test split, real spliced events, real classical
baselines --- and that SENTINEL-E's time-uniform guarantee holds on both
corpora at the lag each one's own diagnostic selects, which the naive
classical alternatives do not manage on either. It does not establish that
this backbone, at this frame rate, is adequate for deployment: Avenue's
power is weak, the horizon tested is minutes rather than weeks, and the two
datasets are pedestrian-walkway anomaly corpora, not illegal-dumping
footage, because no continuously-labelled real dumping corpus of the
required length is publicly available (\cref{sec:limitations}). What
transfers is the pipeline and the diagnostic, not the specific power number:
a stronger backbone with genuinely independent frame-to-frame errors, run
through the identical calibration and episodic layers, is the direction most
likely to close the gap \cref{tab:real_lag_sweep} exposes.
